% Section: Importancia de la protección radiológica
\chapter{Importancia de la Protecci\'on Radiol\'ogica}

El uso de la radiactividad en la ingenier\'ia es uno de los aspectos m\'as difundidos de la tecnolog\'ia nuclear, puesto que forma parte del pequeño número de t\'ecnicas de control de calidad mediante pruebas no destructivas.

Pero as\'i como son muchas las ventajas que nos ofrece tambi\'en se corren muchos riesgos si no sabemos emplearla bajo ciertos procedimientos que garanticen nuestra protecci\'on contra la radiaci\'on.

Estar conscientes de que el programa de protecci\'on radiol\'ogica est\'a hecho para asegurar nuestra es dar el primer paso en la comprensi\'on de este curso.

El objetivo principal de la protecci\'on radiol\'ogica es lograr que las operaciones controladas por la direcci\'on de un centro se desarrollen sin riesgo para la salud y la seguridad de las personas que estén dentro o fuera del \'area, as\'i como la protecci\'on del medio ambiente, controlando los riesgos que se derivan de las actividades que, por las caracter\'isticas de ciertos materiales o equipos que se utilizan, puedan implicar la emisi\'on de radiaciones ionizantes. Para conseguirlo, hay que aplicar los reglamentos que esten en vigor y todas las dem\'as normas pertinentes, y vigilar debidamente las operaciones para comprobar la eficacia de las medidas de protecci\'on adoptadas.

Las responsabilidades que se adquieren al ser nombrado como personal ocupacionalmente expuesto son:
\begin{itemize}
  \item Conocer y aplicar correctamente los principios b\'asicos de seguridad radiol\'ogica.
  \item Conocer el uso correcto de fuentes de radiaci\'on ionizante, equipo detector, y accesorios accesorios.
  \item Usar durante la jornada de trabajo el vestuario, equipo, accesorios y dispositivos que se requerir\'an de acuerdo a lo establecido en el manual de seguridad radiol\'ogica y la norma puesta por el encargado de seguridad radiol\'ogica.
  \item Informar al encargado de seguridad radiol\'ogica sobre cualquier situaci\'on de alto riesgo o incidente radiol\'ogico.
\end{itemize}

% Section: Estructura Atómica
\chapter{Estructura At\'omica}

En la naturaleza, la materia se encuentra en una gran variedad de formas. Desde el punto de vista de la qu\'imica, las sustancias puras se distinguen seg\'un sean compuestos o elementos. Los elementos son sustancias que no pueden descomponerse m\'as mediante métodos qu\'imicos.

La forma m\'inima en que se puede encontrar un compuesto se llama mol\'ecula y la forma m\'inima en que se puede encontrar un elemento se llama \ul{atomo}.
